\chapter[CONCLUSÃO]{Conclusão}
\label{cap:conclusao}

Este trabalho desenvolveu e avaliou um sistema que localiza, na Web aberta,
editais de fomento com inscrições em aberto. O sistema combina a recuperação
de documentos candidatos por meio da API de busca Tavily \cite{tavily} com uma
cadeia de filtros
heurísticos de listagem, relevância e prazo, seguida de deduplicação e
persistência. As heurísticas foram calibradas sobre um conjunto rotulado de 47
itens e avaliadas uma única vez sobre um conjunto independente de 26 itens,
coletado em outra data.

O principal resultado foi a redução dos editais descartados indevidamente, de
5 para 1 entre os 14 editais do conjunto de teste, sem aumento dos itens
errados entregues ao usuário. A extração de prazo, que na configuração original
não identificava corretamente nenhuma data, passou a acertar 4 das 12 datas do
conjunto de teste. Esses ganhos se mantiveram fora da amostra de calibração,
embora menores que os observados nela. A precisão do resultado final, por
outro lado, não melhorou: o sistema continua entregando páginas que não são
editais abertos, sobretudo listagens.

% -----------------------------------------------------------------------------
\section{Contribuições}
\label{sec:contribuicoes}

Além do sistema em si, o trabalho deixa três contribuições de método.

\begin{itemize}
  \item \textbf{Coleta reprodutível}. O registro de cada coleta em um
    \textit{snapshot}, com a data de referência fixada no manifesto, permite
    reprocessar a mesma coleta com configurações diferentes e reproduzir
    exatamente as decisões do sistema, o que é condição para comparar
    configurações em um domínio que muda continuamente.
  \item \textbf{Avaliação orientada ao objetivo do sistema}. As métricas
    operacionais de editais perdidos e de itens errados aceitos, combinadas
    em um critério lexicográfico, evitaram que a calibração otimizasse
    métricas por etapa à custa do resultado entregue ao usuário.
  \item \textbf{Análise de sensibilidade como diagnóstico}. A investigação dos
    parâmetros sem efeito revelou três defeitos de implementação que não eram
    visíveis nas métricas agregadas, entre eles uma extração de prazo cuja
    pontuação de frases nunca atuava. Em sistemas baseados em regras, a
    análise de sensibilidade serve não apenas para ajustar parâmetros, mas
    também para verificar se cada regra de fato participa das decisões.
\end{itemize}

% -----------------------------------------------------------------------------
\section{Limitações e ameaças à validade}
\label{sec:ameacas}

Seguindo a classificação de \citeonline{wohlin2012}, identificaram-se
quatro ameaças à validade dos resultados.

\begin{itemize}
  \item \textbf{Tamanho da amostra}. Com 47 e 26 itens rotulados, e com um e
    três editais abertos, uma diferença de um único item altera as métricas de
    forma considerável. Para mitigar essa ameaça, as decisões apoiaram-se em
    métricas por etapa, a escolha da busca em grade foi validada por validação
    cruzada e as diferenças foram acompanhadas de intervalos de confiança por
    \textit{bootstrap}. Mesmo assim, apenas a redução de editais perdidos tem
    intervalo de confiança que exclui o zero.
  \item \textbf{Sobreajuste ao conjunto de calibração}. Regras como a âncora de
    limite foram concebidas a partir da análise dos erros na calibração, e a
    acurácia da data foi menor no teste do que na calibração. A reserva do
    conjunto de teste, avaliado uma única vez com a configuração congelada,
    permite estimar o desempenho fora da amostra de ajuste \cite{cawley2010},
    mas não elimina o viés das escolhas feitas.
  \item \textbf{Rotulador único}. Os rótulos foram atribuídos por uma só pessoa,
    o que impede medir a concordância entre rotuladores. O guia de rotulação
    escrito e as regras de consistência reduzem, mas não eliminam, a
    subjetividade de rótulos como a relevância.
  \item \textbf{Representatividade das coletas}. As duas coletas usaram o mesmo
    gerador de consultas e foram feitas com onze dias de intervalo. Os
    resultados podem não se estender a outros domínios, a outras épocas do ano
    ou a outras APIs de busca além da Tavily, pois todo o conjunto de
    documentos candidatos dependeu da recuperação feita por ela. A exclusão, no conjunto de teste, dos itens já presentes
    na calibração reduziu a proporção de listagens nesse conjunto. Casos como
    chamadas de fluxo contínuo, editais com vários ciclos de inscrição e itens
    dos quais só se obteve o resumo da busca ficaram sub-representados.
\end{itemize}

Há ainda limitações de projeto. O \textit{pipeline} não tem uma etapa que
decida se a página é um edital, e os itens errados aceitos são, em sua maioria,
listagens que atravessaram todos os filtros. A extração de texto do HTML
descarta a estrutura da página, o que transforma tabelas de cronograma em
frases longas com várias datas e limita a extração de prazo. Por fim, a
geração de consultas e a deduplicação não foram avaliadas quantitativamente:
a calibração concentrou-se nos três filtros heurísticos.

% -----------------------------------------------------------------------------
\section{Trabalhos futuros}
\label{sec:futuros}

As limitações apontam os seguintes desdobramentos:

\begin{itemize}
  \item \textbf{Etapa de identificação de edital}. Incluir no
    \textit{pipeline} uma decisão explícita sobre se a página é um edital,
    cujo rótulo já existe nos conjuntos construídos, para reduzir os itens
    errados aceitos.
  \item \textbf{Extração de texto com estrutura}. Preservar quebras de linha e
    células de tabela na conversão do HTML, de modo que cada linha de um
    cronograma seja tratada como uma frase própria.
  \item \textbf{Ampliação da amostra}. Rotular parte das coletas diárias
    gravadas pelo servidor, para obter mais editais abertos, e incluir um
    segundo rotulador para medir a concordância entre rotuladores. Com esses
    dados, o guia da Seção~\ref{sec:impl_recalibracao} permite conduzir uma
    nova análise sem alterar o código.
  \item \textbf{Comparação com modelos supervisionados}. Com mais dados
    rotulados, comparar as heurísticas calibradas com classificadores
    treinados, nos moldes da comparação de \citeonline{buckmann2017}.
  \item \textbf{Avaliação da recuperação}. Avaliar a geração e a priorização de
    consultas e a deduplicação, comparando estratégias pela proporção de
    editais abertos recuperados.
\end{itemize}
